%%%%%%%%%%  scan idx 118  |  folio 107  %%%%%%%%%%

\medskip
\noindent\textbf{Teorema 4.11.}
Sea $X_0$ un esquema separado de tipo finito sobre $k$, liso fuera
de un subconjunto cerrado finito
\[
X_0^{\mathrm{Sing}}=\{x_1,\ldots,x_r\}.
\]
Sean $U_0^1,\ldots,U_0^r$ entornos abiertos afines de
$x_1,\ldots,x_r$ tales que
$U_0^i\cap X_0^{\mathrm{Sing}}=\{x_i\}$. Si
$H^2(X_0,\underline D^0_{X_0})=0$, entonces el funtor natural
\[
\underline M_{X_0}\longrightarrow
\underline M_{U_0^1}\times\cdots\times\underline M_{U_0^r}
\]
es liso.

\smallskip
\noindent\textit{Demostración.}
Completemos $U_0^1,\ldots,U_0^r$ hasta obtener un recubrimiento
abierto afín de $X_0$, eligiendo $U_0^{r+1},\ldots,U_0^N$ tales que
$U_0^j\cap X_0^{\mathrm{Sing}}=\varnothing$ para todo $j>r$.
Sea $R'\longrightarrow R$ una pequeña extensión en
$\underline C_\Lambda$, y sean
$(R',U'^1),\ldots,(R',U'^r)$ deformaciones de
$U_0^1,\ldots,U_0^r$ sobre $R'$ tales que
\[
U'^i\otimes_{R'}R=X|_{U_0^i},\qquad 1\leq i\leq r,
\]
para alguna deformación $(R,X)$ de $X_0$ sobre $R$. Debemos hallar
una deformación $(R',X')$ de $X_0$ sobre $R'$ tal que
\[
X'\otimes_{R'}R=X,
\qquad
X'|_{U_0^i}=U'^i\quad(1\leq i\leq r).
\]

%%%%%%%%%%  scan idx 119  |  folio 108  %%%%%%%%%%

Para $j>r$, el esquema $U_0^j$ es liso sobre $k$; por tanto, por
4.9(a), $X|_{U_0^j}$ se levanta a una deformación $(R',U'^j)$ de
$U_0^j$ sobre $R'$. Así,
\[
U'^j\otimes_{R'}R=X|_{U_0^j}\qquad(j>r).
\]
Para cada par $i\neq j$, la intersección $U_0^i\cap U_0^j$ es afín
y lisa sobre $k$. Obtenemos, pues, un isomorfismo
\[
\tau_{ij}:U'^i|_{U_0^i\cap U_0^j}
\xrightarrow{\ \sim\ }
U'^j|_{U_0^i\cap U_0^j}
\]
tal que
$\tau_{ij}\otimes_{R'}R=\operatorname{id}_{X|_{U_0^i\cap U_0^j}}$.
En las intersecciones triples, los automorfismos
\[
\tau_{ij}\tau_{jk}\tau_{ik}^{-1}
\]
definen una clase de cohomología en
$H^2(X_0,\underline D^0_{X_0})$. Esta clase es la obstrucción para
pegar los $U'^i$ en una deformación global. Se anula por hipótesis,
de modo que los $U'^i$ se pegan en una deformación $(R',X')$ con
$X'\otimes_{R'}R=X$ y $X'|_{U_0^i}\simeq U'^i$. \hfill$\square$

\medskip
\noindent\textbf{Observación 4.12.}
La demostración precedente supone que $X_0$ es separado, pero esta
hipótesis es innecesaria. En efecto, fijemos una deformación $(R,X)$
de $X_0$ y deformaciones $(R',U'^i)$ de los miembros $U_0^i$ de un
recubrimiento abierto. Las deformaciones $(R',V')$ de los
subconjuntos abiertos $V$ de $X$ sobre $R'$ que son compatibles con
las deformaciones ya dadas $(R',U'^i)$ forman una «gerbe» (en la
terminología de Giraud) cuyo haz estructural es
$\underline D^0_{X_0}$. La obstrucción a la existencia de una sección
de esta gerbe---es decir, de una deformación global $(R',X')$ de $X$
compatible con los datos dados $(R',U'^i)$---pertenece a
$H^2(X_0,\underline D^0_{X_0})$ y, por tanto, se anula bajo la
hipótesis del teorema.
